\section{Application}
\label{sec:application}

This section documents the dual ascent method
used to solve the planner's optimisation problem~\eqref{eq:planning_problem}
at each quarter of the simulation.
It then establishes a systems-optimisation approach for the decentralised firm layer, where a set of local
linear programmes allocate production across competing firms to maximise total income
subject to distributed capital constraints, presents performance results on the 
2022 Spanish input-output benchmark, and establishes a rigorous Two-Agent New Keynesian (TANK) model as a comparative baseline.
It must be noted that the Spanish data is used as a proof of concept with real data, not as a counterfactual policy model.
Spain's 2022 national accounts provide a well-documented, publicly available IO table with sufficient sectoral disaggregation $n = 65$ to constitute a realistic stress test of the planning algorithm's computational and stability properties.

\subsection{Simulation Methodology}
\label{subsec:methodology}


\paragraph{Iterative solution to the Linear Input-Output Problem}

For solving the production system we use an iterative Neumann series approximation instead of a full inversion because the matrix $(I- \bar{A})\\$ of the inversion is $\mathcal{O}(n^3)$ Therefore to reduce the complexity of the algorithm in time and space.

\begin{equation}\label{eq:Iderative inversion}
    \vec{X}^{(k)} = \bar{A}\vec{X}^{(k)} + \vec{F}
\end{equation}
In this approximation the time complexity is reduced to $\mathcal{O}(k \cdot n^2)$ and $\mathcal{O}(n^2)$ respectively.
In the simulation this is done with the function \texttt{neumann\_apply!(  )} defined in the script \texttt{ModelCore.jl}.
The same function is used with the transpose of the $\bar{A}^{\top}$ matrix to calculate the value of the $\vec{\pi}$ from $\vec{\lambda}_{K}$ and $\lambda_{L}$.
The error in this method of approximation is the following: $\rho(A)^{K} \approx \varepsilon $ here $\varepsilon$ represents the error of the $\text{L}^{2}$ norm of $
\| \vec{X}^{*} - \vec{X}^{(K)}\| = \varepsilon \| \vec{X}^{*} \|
$ and $\text{K}$ is the maximum number of iterations used in calculation of the Neumann approximation.

\paragraph{Simulation Methodology Overview}

The simulation embeds a benevolent social planner inside a fully-specified
heterogeneous multi-agent economy.
The implementation is written in Python (orchestration, calibration,
post-processing) with a compiled Julia backend (\texttt{ModelCore.jl})
that performs all hot inner loops.
To minimize cross-language overhead, we employ a \textbf{Vectorized Neumann Expansion}
technique: during calibration, firm-level capital allocations are computed as a
single batch matrix operation in Julia, reducing inter-process communication
latency by $99\%$ compared to sequential firm loops.

\paragraph{Calibration and initialisation.}
The model is calibrated once from the 2022 Spanish National Accounts
input-output table ($n = 65$ sectors), yielding the technology matrix
$A$, the capital-requirement matrix $B$, the labour-coefficient vector
$\vec{l}$, and the initial capital stock $K_0$.
Per-household preference weights $\alpha^h$ are
drawn from a Dirichlet distribution centred on the aggregate benchmark
shares and calibrated so that household income, derived from a stochastic
distribution matrix $\Phi_{h,f}$ over $n_{\text{firms}}$ firms, clears the
initial goods market. In reality this matrix is governed by the relationships of how groups households and firms exchange incomes.

\paragraph{Planner optimisation.}
Given the effective resource bounds $\vec{K}_{\text{eff}}$ and $L_{\text{eff}}$
net of investment pre-commitment, the planner solves the convex programme
in~\eqref{eq:solver problem} via the Nesterov-accelerated dual ascent
described in Sections~\ref{sec:model}--\ref{sec:optimization}.
The solver is implemented natively in Julia as \texttt{solve\_planner\_native},
which operates entirely in log-space to enforce non-negativity of shadow
prices without explicit projection.
Warm-start multipliers from the previous quarter are passed in whenever
available, drastically reducing the iteration count after Q1.
The solver returns the optimal consumption target $\vec{C}^*$, the gross
output target $\vec{X}^*$, the sector-level shadow price vector $\vec{\pi}^*$,
and the marginal utility of income $\mu^*$.

\paragraph{Decentralised firm production.}
Sectoral capital $\vec{K}$ is partitioned across $n_{\text{firms}}$ firms
via a random draw that preserves aggregate sectoral totals.
Each firm solves the local LP~\eqref{eq:firm_lp} natively on the Julia side
using the \textbf{HiGHS Dual Simplex} solver.
By offloading the entire sectoral firm-set as a block-angular constraint
matrix to Julia, the simulation avoids the significant overhead of
Python-based LP interfaces (e.g., SciPy).
This allows for the simultaneous optimization of 250 independent firms
within milliseconds.
The alignment between planner targets ($\vec{X}^*$) and realised
production ($\vec{X}_{\text{actual}}$) is maintained via a
strict $0.1\%$ primal tolerance.

\paragraph{Income distribution and household demand.}
Firm nominal incomes $Y_f = \vec{v}_{\text{eff}}^\top \vec{X}_f$, where
$\vec{v}_{\text{eff}}$ is calibrated from the first-period shadow weights
so that aggregate income equals initial GDP, are distributed to
$n_h$ households via the distribution matrix:
$Y_h = \sum_f \Phi_{h,f} Y_f$.
Each household $h$ maximises a CRRA utility function with preference
weight vector $\alpha^h$ and good-specific curvature parameters
$\sigma_i$, yielding individual demand
$C_{h,i,t} = (w_h\,\alpha_{i,t,h}/(\mu_t P_{i,t}))^{1/\sigma_i}$.
Aggregate household demand is then cleared against the
physical supply $\vec{C}_{\text{actual}} = (I{-}\bar{A})\vec{X}_{\text{actual}}
- \Delta\vec{K} - \vec{G}$.

\paragraph{Capital accumulation and state update.}
Realised gross investment equals the residual of actual output after
netting consumption, government spending, and intermediate demand.
Each firm's capital stock is updated by the standard law of motion:
\begin{equation}
  K_{f,i,t+1} = (1 - \delta) K_{f,i,t} + I_{f,i,t},
\end{equation}
with $\delta = 0.01$ per quarter.
Investment is distributed across firms in proportion to their capital
shares, preserving the block-angular structure of firm balance sheets.
Aggregate capital is recovered as $K_{i,t+1} = \sum_f K_{f,i,t+1}$,
and an accounting identity check $\vec{K} \equiv \sum_f \vec{K}_f$
is enforced after every quarter.

\paragraph{Planner preference update.}
At the close of each quarter the planner updates its aggregate preference
estimate by incorporating the revealed preference signal from the
tâtonnement (equation~\eqref{eq:preference_update}), adding an $\varepsilon$-floor of $10^{-5}$ to prevent
zero-weight sectors, and renormalising to the simplex.
The evolution of the planner's estimate and its distance from the true
(unobserved) household preferences are tracked via the $\ell^2$- and
$\ell^\infty$-norm gaps, which remain bounded throughout the 20-quarter
horizon (Section~\ref{sec:price_dynamics}).

\subsection{Planner's Problem Formulation}

Substituting the growth-augmented technology and investment
pre-commitment of Sections~\ref{sec:model}--\ref{sec:optimization},
the planner's quarterly problem reduces to:
%
\begin{align}\label{eq:solver problem}
\max_{\vec{C}\geq \vec{0}} \quad & \sum_i \alpha^{agg}_i
  \,\frac{C_i^{1-\sigma_i}}{1-\sigma_i} \notag \\
  \quad\text{s.t.}\quad
  & M\vec{C} \leq \vec{K}_{\mathrm{eff}},\quad
  \vec{P}_{L} \cdot \vec{C} \leq L_{\mathrm{eff}},
\end{align}
%
where $M = B(I-\bar{A})^{-1}$ is the capital-requirement matrix
applied directly to consumption,
$\vec{K}_{\mathrm{eff}} = \vec{K}_t - M (\Delta\vec{K}_t + \vec{G}_{t})$ and
$L_{\mathrm{eff}} = L - \vec{P}_{L} \cdot ( \Delta\vec{K}_t + \vec{G}_{t})$
are resource bounds net of investment pre-commitment, and
$\vec{P}_{L} = (I-\bar{A})^{-\top}\vec{l}$ is the
transformed labour coefficient vector precomputed once at calibration.

This is a convex programme: the objective is strictly concave (since
$\sigma_i > 0$ implies $\partial^2/\partial C_i^2 < 0$ for all $i$)
and the feasible set is a convex polytope, so the unique global maximiser exists
whenever the feasible set is non-empty \citep{boyd2004convex}.

\subsection{Dual Formulation and Stationarity}

Rather than solving the primal problem directly, the algorithm operates on the dual.
By associating non-negative Lagrange multipliers (shadow prices) $\vec{\lambda}_K \in \mathbb{R}^n_+$ and $\lambda_L \in \mathbb{R}_+$ with the capital and labour constraints respectively, we form the Lagrangian:
%
\begin{multline} \label{eq:lagrangian effective}
  \mathcal{L}(\vec{C}, \vec{\lambda}_K, \lambda_L) = \hat{W}(\vec{C})
     - \vec{\lambda}_K^\top \bigl(M\vec{C} - \vec{K}_{\mathrm{eff}}\bigr)\\
    - \lambda_L \bigl(\vec{P}_{L} \cdot \vec{C} - L_{\mathrm{eff}}\bigr).
\end{multline}
%
Taking the derivative with respect to consumption $C_i$ and setting to zero yields the closed-form first-order condition (FOC):
%
\begin{multline} \label{eq:foc}
     \alpha^{agg}_i\,C_i^{-\sigma_i} - \bigl(M^\top \vec{\lambda}_K\bigr)_i - \lambda_L P_{L, i} = 0 \\
     \implies \
     C_i = \left(\frac{\alpha^{agg}_i}{\pi_i}\right)^{1/\sigma_i},
\end{multline}
%
where $\pi_i = (M^\top \vec{\lambda}_K)_i + \lambda_L P_{L,i}$ is the
implied final-demand shadow price of good $i$, and correspondingly
$\mu_t = \pi_i / P_{i,t}$ is the marginal utility of income, common
across goods and households at the optimum.

\subsection{Multiplicative Nesterov-Accelerated Dual Ascent}

The dual function $g(\vec{\lambda}_K, \lambda_L) = \max_{\vec{C}} \mathcal{L}(\vec{C}, \vec{\lambda}_K, \lambda_L)$ is minimized over $\vec{\lambda} \geq 0$ to find the optimum shadow prices.
We employ an accelerated dual-ascent scheme using \textbf{Nesterov momentum} applied to the multipliers in log-space. This approach, known as \textbf{Multiplicative Dual Ascent}, ensures that shadow prices remain strictly positive without requiring explicit projection steps.

The solver maintains a state in log-space $\theta_i = \ln \lambda_i$. At each iteration $k$, we evaluate the gradients at an extrapolated point $\vec{y}^{(k)}$, where the log-gradients correspond to the normalized physical slacks:
%
\begin{equation}\label{eq:slacks}
  g_{K,j}^{(k)} = -\frac{s_{K,j}}{K_{\mathrm{eff},j}}, \quad
  g_L^{(k)} = -\frac{s_L}{L_{\mathrm{eff}}},
\end{equation}
%
where $\vec{s}_K = \vec{K}_{\mathrm{eff}} - M\vec{C}(\vec{y})$ and $s_L = L_{\mathrm{eff}} - \vec{P}_{L} \cdot \vec{C}(\vec{y})$.
The shadow prices are updated via a fixed-step gradient descent in log-space:
%
\begin{align}\label{eq:dual_update}
  \ln \lambda_{K,j}^{(k+1)} &= \ln y_{K,j}^{(k)} + \sigma^{(k)} \eta_K g_{K,j}^{(k)} \notag \\
  \ln \lambda_{L}^{(k+1)} &= \ln y_{L}^{(k)} + \sigma^{(k)} \eta_L g_L^{(k)},
\end{align}
%
where the base learning rates $\eta_K \approx 0.4$ and $\eta_L \approx 0.5$ are dynamically scaled by a factor $\sigma^{(k)}$ determined via a \textbf{backtracking line search} on the dual objective value.

To prevent momentum-driven oscillations near the capacity singularities, we implement an \textbf{Adaptive Gradient Restart} rule. The effective iteration counter $k$ is reset to unity ($k = 1$) whenever the extrapolation direction works against the current gradient signal:
%
\begin{equation}
  \langle \ln \vec{y}^{(k)} - \ln \vec{\lambda}^{(k)}, \vec{g}^{(k)} \rangle < 0 \implies \text{Restart}.
\end{equation}
%
Otherwise, the look-ahead point is generated via the standard Nesterov recurrence:
\begin{multline}\label{eq:extrapolation}
  \ln y_{i}^{(k+1)} = \ln \lambda_i^{(k+1)} + \\\left(\frac{k - 1}{k + 2}\right) \bigl(\ln \lambda_i^{(k+1)} - \ln \lambda_i^{(k)}\bigr).
\end{multline}
This formulation allows the solver to handle sectors with vastly different capital intensities unit-invariantly, as the log-space steps represent relative price adjustments (percentage changes) proportional to sectoral over-utilization.


\paragraph{Complexity per Matrix-Vector Operation (MVP).}
Each iteration of~\eqref{eq:dual_update} requires at least \textbf{two Neumann Applications}: one forward operator call $\mathcal{M}(\vec{C})$ and one adjoint call $\mathcal{M}^\top(\vec{\lambda})$. Each Neumann application of order $k$ corresponds to $k$ atomic \textbf{Matrix-Vector Products (MVPs)} of the technology matrix $A$. The total computational effort per quarter is tracked via cumulative Neumann calls (typically \textbf{100--250}) and total atomic MVPs (typically \textbf{2,000--5,000} for $k=20$).
\subsection{Convergence Criteria}

The iterative loop monitors the physical constraint violations and checks for mathematical stabilization.

\paragraph{KKT Tolerance.}
The algorithm terminates successfully if both the dual complementarity and primal feasibility conditions are satisfied at a high-precision threshold (typically $10^{-4}$ or $10^{-5}$):
\begin{align}
  \text{Dual: } & \lambda_{L}^{(t)} |s_L| \leq \text{tol}_d, \quad \lambda_{K,j}^{(t)}|s_{K,j}| \leq \text{tol}_d \\
  \text{Primal: } & \frac{|s_L|}{L_{\text{eff}}} \leq \text{tol}_p, \quad \frac{| s_{K,j}|}{K_{\text{eff},j}} \leq \text{tol}_p
\end{align}
where $s_L$ and $\vec{s}_K$ are the slacks net of resource bounds.

\subsection{Empirical Complexity Study of the Optimisation Kernel}
\label{subsec:scaling}

To verify the theoretical complexity of the dual-ascent solver independently of
the full simulation stack, we provide a standalone Julia script \texttt{scaling.jl}.
This script is a purpose-built empirical scaling study of
\texttt{solve\_planner\_general}---a dimensionally unrestricted replica of the
production solver \texttt{solve\_planner\_native} in \texttt{ModelCore.jl}---and
has no dependency on the simulation's calibration data or agent layer.

\paragraph{Experimental design.}
The script generates a sequence of synthetic central-planning problems with
problem dimension $n \in \{1{,}000,\, 2{,}000,\, \ldots,\, 10{,}000\}$.
For each $n$, it constructs a pair of $n \times n$ sparse matrices: a Leontief
technology matrix $A$ (spectral radius $\approx 0.15$) and a capital-requirements
matrix $B$, both at density $7.5\,\%$, well below the $10\,\%$ threshold at which
sparse arithmetic retains its asymptotic advantage.  All remaining inputs
(preference weights $\boldsymbol{\alpha}$, curvature parameters $\boldsymbol{\sigma}$,
labour coefficients $\boldsymbol{P_L}$,
capital supply $\boldsymbol{K}$, government demand $\boldsymbol{G}$, and
warm-start duals) are drawn from calibrated random distributions that reflect the
structure of realistic IO data.  A single JIT warm-up pass on a $50$-sector toy
problem is executed before the timed loop so that Julia compilation costs are
excluded from all measurements.

\paragraph{Measured quantities.}
For each $n$ the script records:
\begin{itemize}
  \item \textbf{Iteration count} --- the number of Nesterov AGD steps required
        to satisfy the KKT tolerance $\varepsilon = 10^{-4}$, subject to a
        ceiling of $2{,}000$ iterations.
  \item \textbf{Wall-clock time} --- the elapsed time (seconds) from the first
        gradient evaluation to solver return, measured via Julia's \texttt{time()}.
\end{itemize}
Both quantities are saved as a combined figure:
\texttt{complexity\_scaling\_fit.png} shows wall-clock time with a high-precision 
quadratic fit and the corresponding regression equation.

\paragraph{Complexity findings.}
A quadratic polynomial $t(n) = a n^2 + b n + c$ is fitted by
ordinary least squares to the wall-clock measurements.
The resulting regression found $a \approx 9.20 \times 10^{-7}$, $b \approx -3.59 \times 10^{-3}$, and $c \approx 4.82$.
The dominant quadratic term confirms that the total arithmetic complexity of the
CRRA solver kernel grows as $\mathcal{O}(n^2)$ for the constant-density sparse problem
instances ($d = 15\%$) examined. This matches the theoretical cost of the
sparse matrix--vector products ($O(d n^2)$) and indicates that the Nesterov iteration count remains roughly stable as the problem scales, 
ensuring the system remains performant for large-scale economic planning.

\subsection{Decentralised Firm Production and Household Income Distribution}

To model a distributed economy, production is not computed as a single macroeconomic aggregate. Instead, sectoral capital is partitioned across $f = 1 \dots n_{\text{firms}}$ independent firms. Each firm is endowed with private capital $\vec{K}_f$ and a proportional labour endowment $L_f$.

Due to the linear nature of the value-added objective, the classical optimisation of $\vec{v}_{\text{eff}}^\top \vec{x}_f$ exhibits non-unique optima on a flat optimal face. To mathematically emulate an \textit{Optimum Hopping} mechanism---where firms select the unique vertex combination that maximally aligns with the planner's physical output targets $\vec{X}^*$---we construct a Lexicographic Linear Programme. The tie-breaker is implemented via an $L_1$ penalty using auxiliary variables $\vec{u}_f$:

\begin{align}\label{eq:firm_lp}
\max_{\vec{x}_f \geq 0,\, \vec{u}_f \ge 0} &\quad \vec{v}_{\text{eff}}^\top \vec{x}_f - \epsilon \sum_i u_{f,i} \\
\text{s.t.} &\quad B \vec{x}_f \leq \vec{K}_f \notag\\
            &\quad u_{f,i} \geq \| x_{f,i} - (\text{share}_f \cdot X^*_i) \| \notag\\
\end{align}

Here, $\epsilon$ is chosen sufficiently small (e.g., $10^{-3}$) such that it acts strictly as a tie-breaker on the optimal face. By penalising the Manhattan distance to the proportional target, the solver algorithmically finds the unique convex combination of optimal vertices that spans $\vec{X}^*$.

The total nominal income generated by these firms, $Y_{\text{firm}} = \sum_f \vec{v}_{\text{eff}}^\top \vec{x}_f$, is then distributed heterogeneously to $h = 1 \dots n_h$ households. This distribution is governed by a stochastic ownership matrix $\Phi_{h,f}$, such that each household's income is derived from its dividend share of firm revenues: $Y_h = \sum_f \Phi_{h,f} Y_f$. Consumer demand is then computed at the individual household level using the CRRA demand system~\eqref{eq:demand},
$C_{h,i,t} = (w_h\,\alpha_{i,t,h}/(\mu_t P_{i,t}))^{1/\sigma_i}$,
and dynamically clears against the actual production achievable by the firms $\vec{C}_{\text{actual}} = (I-\bar{A})\vec{X}_{\text{actual}} - (\Delta \vec{K} + \vec{G})$ during the monthly price tâtonnement.

\subsection{Benchmarking Against Decentralised Markets: The HANK Reference}
\label{subsec:hank_benchmark}

To evaluate the efficiency of the cybernetic planner, we contrast it with a high-fidelity \textbf{Heterogeneous Agent New Keynesian (HANK)} model operating on the same 65-sector Spanish input-output structure. In this benchmark, the terminal supply constraints enforced by the planner are replaced by a feedback-driven market coordination mechanism governed by independent households, price-setting firms, and a central bank.

\paragraph{Household Heterogeneity and Liquidity.}
The HANK benchmark simulates $1{,}000$ independent households experiencing idiosyncratic labor income risk (AR(1) process in logs, $\rho = 0.97$). Households are initialized with a Spanish-calibrated wealth distribution ($\text{Gini} \approx 0.67$) and face a hard borrowing constraint ($a \geq 0$). This creates a realistic distribution of \textbf{Marginal Propensities to Consume (MPC)}: households near the liquidity floor exhibit an MPC of $1.0$ (effectively Hand-to-Mouth), while wealthier households smooth consumption via an Euler equation.

\paragraph{NK Pricing and Monetary Policy.}
Each of the 65 sectors is modelled as a set of monopolistically competitive firms setting prices via a log-linearised \textbf{Rotemberg Phillips Curve}. Sectoral inflation $\pi_{i,t}$ is driven by real marginal costs---comprising both labor and intermediate input requirements from the $A$ matrix---relative to a steady-state markup. The nominal interest rate $i_t$ is set by a central bank following an inertial Taylor rule:
\begin{equation}
  i_t = \rho_i i_{t-1} + (1 - \rho_i) \bigl[ r_{ss} + \phi_{\pi} (\pi_t - \pi^*) + \phi_y (y_t - y^*) \bigr],
\end{equation}
where $\pi_t$ is aggregate Laspeyres inflation and $y_t - y^*$ is the output gap relative to potential growth. The real rate $r_t$ is derived via the Fisher approximation.


\paragraph{Uncoordinated Capital Feedback.}
While the planner strictly enforces $K \ge BX$ as a terminal constraint, the TANK benchmark models investment as a reactive feedback loop. Firms demand new capital $\Delta K$ based on current output utilization signals. This creates a reflexive multiplier: higher consumption demands higher investment, which in turn raises nominal income and further stimulates consumption. This "soft" capital adjustment often leads to higher raw GDP growth ($15.5\%$ vs $12.3\%$) at the cost of significant allocative inefficiency, as sectors over-produce goods that are mismatched with the true (unobserved) preference weights $\alpha^h$, ultimately yielding an $8.3\%$ welfare penalty relative to the cybernetic planner.
